\textbf{Phần 1. Trắc nghiệm lựa chọn}
\setcounterpageref{ex}{0}

\begin{ex} %[1D1V04]
Trong các số: $8, 3, 10, 6, 9$, số lớn nhất là số nào?
\choice
{$0$}
{$20$}
{$9$}
{\True $10$}
\loigiai{
So sánh các số đã cho: $3 < 6 < 8 < 9 < 10$.\nSố lớn nhất trong các số trên là số $10$.
\textbf{Đáp số:} $10$.\par
\textbf{Đáp án của câu hỏi là D}
}
\end{ex}

\begin{ex} %[1D1V04]
Trong các số: $7, 2, 5, 0, 9$, số bé nhất là số nào?
\choice
{$1$}
{\True $0$}
{$2$}
{$3$}
\loigiai{
Số bé nhất trong dãy số tự nhiên từ 0 đến 10 là số $0$.
\textbf{Đáp số:} $0$.\par
\textbf{Đáp án của câu hỏi là B}
}
\end{ex}

\begin{ex} %[1D1V04]
Tìm số tự nhiên $x$ thỏa mãn: $6 < x < 8$.
\choice
{$6$}
{$8$}
{\True $7$}
{$5$}
\loigiai{
Số tự nhiên nằm giữa 6 và 8 là số liền sau của 6 và liền trước của 8, đó là số $7$.
\textbf{Đáp số:} $7$.\par
\textbf{Đáp án của câu hỏi là C}
}
\end{ex}

\begin{ex} %[1D1V04]
Có bao nhiêu số tự nhiên $x$ thỏa mãn: $4 < x < 9$?
\choice
{\True $4$}
{$3$}
{$5$}
{$2$}
\loigiai{
Các số tự nhiên lớn hơn 4 và bé hơn 9 là: $5, 6, 7, 8$.\nĐếm các số trên có tất cả $4$ số.
\textbf{Đáp số:} $4$.\par
\textbf{Đáp án của câu hỏi là A}
}
\end{ex}

\begin{ex} %[1D1V04]
Cho các số: $9, 4, 7, 2, 10$. Sắp xếp các số theo thứ tự từ bé đến lớn, số đứng ở vị trí thứ ba là số nào?
\choice
{$6$}
{$8$}
{$5$}
{\True $7$}
\loigiai{
Sắp xếp các số từ bé đến lớn: $2 < 4 < 7 < 9 < 10$.\nSố đứng ở vị trí thứ ba là số $7$.
\textbf{Đáp số:} $7$.\par
\textbf{Đáp án của câu hỏi là D}
}
\end{ex}

\begin{ex} %[1D1V04]
Cho các số: $10, 5, 8, 3, 7$. Sắp xếp các số theo thứ tự từ lớn đến bé, số đứng ở vị trí thứ hai là số nào?
\choice
{$7$}
{\True $8$}
{$9$}
{$6$}
\loigiai{
Sắp xếp các số từ lớn đến bé: $10 > 8 > 7 > 5 > 3$.\nSố đứng ở vị trí thứ hai là số $8$.
\textbf{Đáp số:} $8$.\par
\textbf{Đáp án của câu hỏi là B}
}
\end{ex}

\begin{ex} %[1D1V04]
Tìm số lớn nhất có một chữ số bé hơn 9.
\choice
{$7$}
{$9$}
{\True $8$}
{$6$}
\loigiai{
Các số có một chữ số bé hơn 9 là: $0, 1, 2, 3, 4, 5, 6, 7, 8$.\nSố lớn nhất trong các số này là số $8$.
\textbf{Đáp số:} $8$.\par
\textbf{Đáp án của câu hỏi là C}
}
\end{ex}

\begin{ex} %[1D1V04]
Tìm số bé nhất có một chữ số lớn hơn 5.
\choice
{\True $6$}
{$5$}
{$7$}
{$4$}
\loigiai{
Các số có một chữ số lớn hơn 5 là: $6, 7, 8, 9$.\nSố bé nhất trong các số này là số $6$.
\textbf{Đáp số:} $6$.\par
\textbf{Đáp án của câu hỏi là A}
}
\end{ex}

\begin{ex} %[1D1V04]
Lấy số lớn nhất trong các số $6, 2, 10, 7$ trừ đi số bé nhất trong các số đó, ta được kết quả bằng bao nhiêu?
\choice
{$7$}
{$9$}
{$6$}
{\True $8$}
\loigiai{
Trong các số đã cho:\n- Số lớn nhất là số $10$.\n- Số bé nhất là số $2$.\nHiệu của số lớn nhất và số bé nhất là: $10 - 2 = 8$.
\textbf{Đáp số:} $8$.\par
\textbf{Đáp án của câu hỏi là D}
}
\end{ex}

\begin{ex} %[1D1V04]
Có bao nhiêu số tự nhiên lớn hơn 2 và bé hơn hoặc bằng 8?
\choice
{$5$}
{\True $6$}
{$7$}
{$4$}
\loigiai{
Các số tự nhiên đó là: $3, 4, 5, 6, 7, 8$.\nĐếm các số trên có tất cả $6$ số.
\textbf{Đáp số:} $6$.\par
\textbf{Đáp án của câu hỏi là B}
}
\end{ex}

\begin{ex} %[1D1V04]
Tìm số tự nhiên $a$ sao cho: $a + 1 = 9$. Giá trị của $a$ là bao nhiêu?
\choice
{$7$}
{$9$}
{\True $8$}
{$6$}
\loigiai{
Muốn tìm số hạng $a$, ta lấy tổng trừ đi 1: $a = 9 - 1 = 8$.
\textbf{Đáp số:} $8$.\par
\textbf{Đáp án của câu hỏi là C}
}
\end{ex}

\begin{ex} %[1D1V04]
Tìm số tự nhiên $b$ sao cho: $10 - b = 3$. Giá trị của $b$ là bao nhiêu?
\choice
{\True $7$}
{$6$}
{$8$}
{$5$}
\loigiai{
Muốn tìm số trừ $b$, ta lấy số bị trừ là 10 trừ đi hiệu là 3: $b = 10 - 3 = 7$.
\textbf{Đáp số:} $7$.\par
\textbf{Đáp án của câu hỏi là A}
}
\end{ex}

\textbf{Phần 2. Trắc nghiệm đúng sai}
\setcounterpageref{ex}{0}

\begin{ex} %[1D1V04]
Xét tính đúng hoặc sai của các khẳng định toán học sau đây:
\choiceTF
{Số thích hợp cần điền vào $5 < \dots < 7$ là $7$.}
{Kết quả của bài toán «Cho ba số $a, b, c$. Biết $a > b$ và $b > c$. Nếu $a = 9$ và $c = 7$, hỏi $b$ bằng bao nhiêu» là $9$.}
{\True Kết quả của bài toán «Tìm số tự nhiên $x$ lớn nhất thỏa mãn: $x < 10$» là $9$.}
{\True Kết quả của bài toán «Tìm số tự nhiên $y$ bé nhất thỏa mãn: $y > 6$» là $7$.}
\loigiai{
\begin{itemize}
    \item Ý a) Sai: Kết quả đúng là $6$, không phải $7$. Số tự nhiên đứng giữa 5 và 7 là số $6$.
\textbf{Đáp số:} $6$.
    \item Ý b) Sai: Kết quả đúng là $8$, không phải $9$. Theo đề bài: $7 < b < 9$.\nSố tự nhiên nằm giữa 7 và 9 là số $8$. Vậy $b = 8$.
\textbf{Đáp số:} $8$.
    \item Ý c) Đúng: Các số tự nhiên bé hơn 10 là: $0, 1, 2, \dots, 9$.\nSố lớn nhất trong các số này là số $9$.
\textbf{Đáp số:} $9$.
    \item Ý d) Đúng: Các số tự nhiên lớn hơn 6 là: $7, 8, 9, 10, \dots$\nSố bé nhất trong các số này là số $7$.
\textbf{Đáp số:} $7$.
\end{itemize}
\begin{itemize}
    \item a) \textbf{Sai}.
    \item b) \textbf{Sai}.
    \item c) \textbf{Đúng}.
    \item d) \textbf{Đúng}.
\end{itemize}
}
\end{ex}

\begin{ex} %[1D1V04]
Xét tính đúng hoặc sai của các bài toán và phép tính sau đây:
\choiceTF
{\True Kết quả của bài toán «Có bao nhiêu số tự nhiên có một chữ số lớn hơn 3» là $6$.}
{Kết quả của bài toán «Tổng của số lớn nhất có một chữ số và số bé nhất có một chữ số là bao nhiêu» là $8$.}
{Kết quả của bài toán «Hiệu giữa số bé nhất có hai chữ số và số lớn nhất có một chữ số bằng bao nhiêu» là $0$.}
{\True Kết quả của bài toán «Tìm số tự nhiên $x$ lớn nhất thỏa mãn: $x < 10$» là $9$.}
\loigiai{
\begin{itemize}
    \item Ý a) Đúng: Các số có một chữ số lớn hơn 3 là: $4, 5, 6, 7, 8, 9$.\nĐếm các số này có tất cả $6$ số.
\textbf{Đáp số:} $6$.
    \item Ý b) Sai: Kết quả đúng là $9$, không phải $8$. Số lớn nhất có một chữ số là số $9$.\nSố bé nhất có một chữ số là số $0$.\nTổng của hai số là: $9 + 0 = 9$.
\textbf{Đáp số:} $9$.
    \item Ý c) Sai: Kết quả đúng là $1$, không phải $0$. Số bé nhất có hai chữ số là số $10$.\nSố lớn nhất có một chữ số là số $9$.\nHiệu của hai số là: $10 - 9 = 1$.
\textbf{Đáp số:} $1$.
    \item Ý d) Đúng: Các số tự nhiên bé hơn 10 là: $0, 1, 2, \dots, 9$.\nSố lớn nhất trong các số này là số $9$.
\textbf{Đáp số:} $9$.
\end{itemize}
\begin{itemize}
    \item a) \textbf{Đúng}.
    \item b) \textbf{Sai}.
    \item c) \textbf{Sai}.
    \item d) \textbf{Đúng}.
\end{itemize}
}
\end{ex}

\textbf{Phần 3. Trả lời ngắn}
\setcounterpageref{ex}{0}

\begin{ex} %[1D1V04]
Cho dãy số: $2, 4, 6, 8, 10$. Số đứng ở vị trí thứ tư trong dãy là số nào?
\shortans{8}
\loigiai{
\textbf{Đáp số:} $8$\par
Đếm theo thứ tự từ trái sang phải:\n- Thứ nhất là 2; thứ hai là 4; thứ ba là 6; thứ tư là 8; thứ năm là 10.\nVậy số đứng ở vị trí thứ tư là số $8$.
\textbf{Đáp số:} $8$.
}
\end{ex}

\begin{ex} %[1D1V04]
So sánh hai phép tính: $4 + 5$ và $3 + 6$. Kết quả của mỗi phép tính đều bằng bao nhiêu?
\shortans{9}
\loigiai{
\textbf{Đáp số:} $9$\par
Ta có: $4 + 5 = 9$ và $3 + 6 = 9$.\nKết quả của hai phép tính đều bằng $9$.
\textbf{Đáp số:} $9$.
}
\end{ex}

\begin{ex} %[1D1V04]
Số nào trong các số $6, 7, 8, 9$ khi cộng với 2 thì được kết quả lớn hơn 9?
\shortans{8}
\loigiai{
\textbf{Đáp số:} $8$\par
Thử từng số:\n- $6 + 2 = 8 < 9$.\n- $7 + 2 = 9$ (bằng 9).\n- $8 + 2 = 10 > 9$ (chọn).\nSố cần tìm là số $8$.
\textbf{Đáp số:} $8$.
}
\end{ex}

\begin{ex} %[1D1V04]
Tìm số tự nhiên $m$ biết rằng: $8 - m > 5$ và $m$ là số tự nhiên lớn nhất có thể. Giá trị của $m$ là bao nhiêu?
\shortans{2}
\loigiai{
\textbf{Đáp số:} $2$\par
Ta có: $8 - 1 = 7 > 5$; $8 - 2 = 6 > 5$; $8 - 3 = 5$ (không lớn hơn 5).\nVậy các số $m$ thỏa mãn là 0, 1, 2. Giá trị lớn nhất của $m$ là số $2$.
\textbf{Đáp số:} $2$.
}
\end{ex}

\begin{ex} %[1D1V04]
Trong một cuộc thi chạy, bạn Nam về đích sau bạn An nhưng trước bạn Bình. Biết bạn An về đích thứ nhất, bạn Bình về đích thứ 3. Hỏi bạn Nam về đích thứ mấy?
\shortans{2}
\loigiai{
\textbf{Đáp số:} $2$\par
Bạn Nam về đích sau người thứ nhất và trước người thứ ba, vậy bạn Nam về đích ở vị trí thứ $2$.
\textbf{Đáp số:} $2$.
}
\end{ex}

\begin{ex} %[1D1V04]
Hiệu của số đứng liền sau số 8 và số đứng liền trước số 4 là bao nhiêu?
\shortans{6}
\loigiai{
\textbf{Đáp số:} $6$\par
Số đứng liền sau số 8 là số $9$.\nSố đứng liền trước số 4 là số $3$.\nHiệu của hai số là: $9 - 3 = 6$.
\textbf{Đáp số:} $6$.
}
\end{ex}
